\section*{Bab \ref{ch:g6:inside-the-earth} --- Bagian Dalam Bumi}

\begin{solution}{exo:g6:inside-the-earth:1}
\omterm{def:g6:inside-the-earth:layers}{Kerak}, \omterm{def:g6:inside-the-earth:layers}{mantel}, \omterm{def:g6:inside-the-earth:layers}{inti luar}, \omterm{def:g6:inside-the-earth:layers}{inti dalam}. \omterm{def:g6:inside-the-earth:layers}{Inti luar} itu cair. \omterm{def:g6:inside-the-earth:layers}{Inti luar} dan
\omterm{def:g6:inside-the-earth:layers}{inti dalam} tersusun dari logam (terutama besi, dengan sedikit nikel); \omterm{def:g6:inside-the-earth:layers}{kerak}
dan \omterm{def:g6:inside-the-earth:layers}{mantel} tersusun dari \omterm{def:g3:stones-pebbles-sand:rock}{batuan}.
\end{solution}

\begin{solution}{exo:g6:inside-the-earth:2}
(a) Salah: \omterm{def:g6:inside-the-earth:layers}{kerak} lebih tebal di bawah \omterm{def:g3:continents-and-oceans:continent}{benua} (kira-kira 30~km) daripada di
bawah \omterm{def:g2:water-on-earth:ocean}{samudra} (kira-kira 5~km). (b) Salah: \omterm{def:g6:inside-the-earth:layers}{inti} tersusun dari logam, paduan
besi dan nikel. (c) Benar. (d) Benar: 12\,262~m, kira-kira 12.3~km.
\end{solution}

\begin{solution}{exo:g6:inside-the-earth:3}
(a) Pada 2891~km, kira-kira 2900~km. (b) Pada 5150~km. (c) $5150 - 2891 = 2259$:
kira-kira 2260~km.
\end{solution}

\begin{solution}{exo:g6:inside-the-earth:4}
``Kulit telur seperti \underline{\omterm{def:g6:inside-the-earth:layers}{kerak}}, putih telur seperti
\underline{\omterm{def:g6:inside-the-earth:layers}{mantel}} dan kuning telur seperti \underline{\omterm{def:g6:inside-the-earth:layers}{inti}}. Akan tetapi,
berbeda dengan kuning telur, \underline{\omterm{def:g6:inside-the-earth:layers}{inti}} terdiri atas dua bagian:
\underline{\omterm{def:g6:inside-the-earth:layers}{inti luar}} yang cair dan \underline{\omterm{def:g6:inside-the-earth:layers}{inti dalam}} yang padat.''
\end{solution}

\begin{solution}{exo:g6:inside-the-earth:5}
$8.849 \div 6371 \approx 0.00139$, dan $0.00139 \times 100 \approx 0.14$:
tinggi Everest hanya kira-kira \textbf{0.14~\%} jari-jari Bumi. Dilihat
dari jauh, Bumi adalah bola yang sangat mulus: \omterm{def:g1:landscapes-around-us:landform}{gunung} tertingginya hanyalah
tonjolan yang sangat kecil.
\end{solution}

\begin{solution}{exo:g6:inside-the-earth:6}
(a) $30 \div 5 = 6$: enam kali lebih tebal. (b) $100 \div 5 = 20$: dua puluh
kali lebih tebal.
\end{solution}

\begin{solution}{exo:g6:inside-the-earth:7}
(a) $12 + 1 \times 25 = 37~^\circ$C. (b) $12 + 1.5 \times 25 = 12 + 37.5 =
49.5~^\circ$C. (c) $12 + 2 \times 25 = 62~^\circ$C.
\end{solution}

\begin{solution}{exo:g6:inside-the-earth:8}
Kenaikannya $74 - 14 = 60$ derajat sepanjang 2~km, jadi $60 \div 2 =
\textbf{30}$~derajat per kilometer. Ya: 30 terletak di antara rata-rata
bagian timur (25) dan rata-rata bagian barat (34).
\end{solution}

\begin{solution}{exo:g6:inside-the-earth:9}
(a) Tepat di atas \omterm{def:g6:inside-the-earth:layers}{inti}, kira-kira 13.7~km per detik; tepat di bawah bagian
atasnya, kira-kira 8.1~km per detik: gelombang P tiba-tiba melambat.
(b) Di antara 2891~km dan 5150~km, di \omterm{def:g6:inside-the-earth:layers}{inti luar}. (c) Gelombang S tidak dapat
merambat melalui zat cair: \omterm{def:g6:inside-the-earth:layers}{inti luar} itu cair.
\end{solution}

\begin{solution}{exo:g6:inside-the-earth:10}
(a) $12.3 \div 30 = 0.41$: kira-kira \textbf{41~\%} dari \omterm{def:g6:inside-the-earth:layers}{kerak} \omterm{def:g3:continents-and-oceans:continent}{benua} yang
rata-rata. (b) Tidak: di bawah \omterm{def:g3:continents-and-oceans:continent}{benua}, \omterm{def:g6:inside-the-earth:layers}{mantel} dimulai di bawah \omterm{def:g6:inside-the-earth:layers}{kerak},
rata-rata kira-kira 30~km di bawah tanah; lubang bor itu berhenti jauh di
atasnya, di dalam \omterm{def:g6:inside-the-earth:layers}{kerak}.
\end{solution}

\begin{solution}{exo:g6:inside-the-earth:11}
(a) $29 \times 6371 = 184\,759$, dan $15 + 184\,759 = 184\,774$: kira-kira
185\,000~derajat. (b) Laboratorium memberikan kira-kira 6000~derajat di
tepi \omterm{def:g6:inside-the-earth:layers}{inti dalam}: angka Leo kira-kira 30~kali terlalu besar
($184\,774 \div 6000 \approx 31$). (c) Laju 29~derajat per kilometer hanya
diukur pada kira-kira 2~km pertama. Leo memakainya ribuan kilometer di luar
jangkauan pengukuran; jauh di bawah, suhu naik jauh lebih lambat. Sebuah
laju hanya boleh dipakai di tempat laju itu diukur.
\end{solution}

\begin{solution}{exo:g6:inside-the-earth:12}
(a) Jika \omterm{def:g6:inside-the-earth:layers}{kerak} lebih ringan daripada rata-ratanya, bahan yang lebih dalam
pasti mempunyai massa jenis lebih dari 5.5, agar rata-ratanya naik.
(b) Di \omterm{def:g6:inside-the-earth:layers}{inti luar} dan \omterm{def:g6:inside-the-earth:layers}{inti dalam} (dari kira-kira 9.9 sampai kira-kira
13~gram per sentimeter kubik); \omterm{def:g6:inside-the-earth:layers}{mantel} hanya mencapai 5.6 tepat di dasarnya.
(c) Bahan yang massa jenisnya lebih besar tenggelam di bawah bahan yang
lebih ringan, seperti batu yang tenggelam di dalam air: logam yang berat
berakhir di tengah dan \omterm{def:g3:stones-pebbles-sand:rock}{batuan} yang lebih ringan di atas.
\end{solution}

\begin{solution}{pb:g6:inside-the-earth:1}
\textbf{1.} $2 \times 6371 = 12\,742$~km. Model itu selebar 2~m, yaitu
2000~mm.

\textbf{2.} $12\,742 \div 2000 \approx 6.37$: 1~mm pada model mewakili
kira-kira 6.37~km Bumi yang sebenarnya.

\textbf{3.} $8.849 \div 6.371 \approx 1.39$: tinggi Everest kira-kira
\textbf{1.4~mm} pada model, sebesar sebutir \omterm{def:g3:stones-pebbles-sand:sand}{pasir}.

\textbf{4.} $12.3 \div 6.371 \approx 1.9$~mm. Pusat model terletak 1000~mm
(1~m) di bawah permukaan: lubang bor itu kira-kira 520~kali lebih pendek
($1000 \div 1.9 \approx 520$).

\textbf{5.} $2891 \div 6.371 \approx 453.8$: bagian atas \omterm{def:g6:inside-the-earth:layers}{inti luar} dicat
kira-kira \textbf{454~mm} di bawah permukaan.

\textbf{6.} $5150 \div 6.371 \approx 808.3$: kira-kira 808~mm di bawah
permukaan. Jari-jari \omterm{def:g6:inside-the-earth:layers}{inti dalam} pada model adalah $1000 - 808 = 192$~mm.

\textbf{7.} \omterm{def:g6:inside-the-earth:layers}{Mantel}: dari kira-kira 5~mm sampai 454~mm, kira-kira 449~mm. \omterm{def:g6:inside-the-earth:layers}{Inti
luar} mengambil $808 - 454 = 354$~mm dan \omterm{def:g6:inside-the-earth:layers}{inti dalam} 192~mm.

\textbf{8.} $30 \div 6.371 \approx 4.7$~mm: sebuah pita tipis, yang dapat
dicat dengan kuas halus.

\textbf{9.} $15 + 2 \times 29 = 73~^\circ$C. Pada model, 2~km terletak
$2 \div 6.371 \approx 0.31$~mm di bawah permukaan: hampir tidak lebih dari
tebal tiga lembar kertas.

\textbf{10.} $4 \div 6.371 \approx 0.63$~mm, dibandingkan dengan kira-kira
1.9~mm untuk lubang bor: tambang itu lebih dangkal, kira-kira
$1.9 - 0.6 = 1.3$~mm.

\textbf{11.} Meleleh atau tidaknya suatu zat bergantung pada suhu dan pada
tekanan. \omterm{def:g6:inside-the-earth:layers}{Inti dalam} tertekan paling kuat, kira-kira 3.3~juta kali tekanan
udara: besi yang tertekan baru meleleh pada suhu yang lebih tinggi, jadi
besi di \omterm{def:g6:inside-the-earth:layers}{inti dalam} tetap padat walaupun lebih panas daripada \omterm{def:g6:inside-the-earth:layers}{inti luar}.

\textbf{12.} $5 \div 6.371 \approx 0.785$: \omterm{def:g6:inside-the-earth:layers}{kerak} \omterm{def:g2:water-on-earth:ocean}{samudra} hanya boleh dicat
setebal \textbf{0.78~mm} pada model selebar 2 meter, kira-kira setebal
delapan lembar kertas ($0.78 \div 0.1 \approx 8$).
\end{solution}
